% BTS CIEL 1 — Réseaux — 09 — FTP
% Compiler deux fois avec XeLaTeX depuis fichiers/.
\documentclass[aspectratio=169,11pt]{beamer}
\usepackage{fontspec}
\setmainfont{Latin Modern Sans}
\setsansfont{Latin Modern Sans}
\setmonofont{Latin Modern Mono}
\usepackage{polyglossia}
\setmainlanguage{french}
\usepackage{tikz}
\usetikzlibrary{arrows.meta,positioning,calc,shapes.geometric}
\usepackage{listings}
\usepackage{xcolor}
\definecolor{foret}{HTML}{1C322D}
\definecolor{creme}{HTML}{F8F3EE}
\definecolor{or}{HTML}{EBB552}
\definecolor{rose}{HTML}{F1CDBE}
\definecolor{vertclair}{HTML}{DDE9DF}
\definecolor{encre}{HTML}{24372E}
\setbeamercolor{normal text}{fg=creme,bg=foret}
\setbeamercolor{frametitle}{fg=or,bg=foret}
\setbeamercolor{title}{fg=creme}
\setbeamercolor{block title}{fg=foret,bg=or}
\setbeamercolor{block body}{fg=encre,bg=creme}
\setbeamercolor{alerted text}{fg=or}
\setbeamercolor{structure}{fg=or}
\setbeamertemplate{navigation symbols}{}
\setbeamertemplate{footline}{\hfill\color{or}\insertframenumber\hspace{0.5cm}\vspace{0.25cm}}
\setbeamertemplate{frametitle}{\vspace{0.35cm}\insertframetitle\par\vspace{0.15cm}\color{or}\rule{\linewidth}{0.6pt}}
\setlength{\parskip}{5pt}
\lstdefinestyle{webcode}{basicstyle=\ttfamily\small\color{encre},keywordstyle=\color{foret},stringstyle=\color{foret},
  backgroundcolor=\color{creme},frame=single,rulecolor=\color{or},framerule=1pt,
  xleftmargin=5pt,xrightmargin=5pt,columns=fullflexible,keepspaces=true,showstringspaces=false}
\newcommand{\question}[1]{\begin{beamercolorbox}[rounded=true,sep=12pt]{block body}\Large\textbf{#1}\end{beamercolorbox}}
\newcommand{\etiquette}[1]{{\color{or}\small\MakeUppercase{#1}}}
% Les schémas ont une surface claire, également lisible en noir et blanc.
\tikzset{
  boite/.style={draw=encre,fill=creme,text=encre,rounded corners,
    align=center,text width=2.8cm,minimum height=1cm},
  flux/.style={-{Stealth},thick,draw=encre},
  every node/.style={text=encre,font=\small}}
\newenvironment{schema}{%
  \begin{beamercolorbox}[sep=8pt,rounded=true]{block body}
  \centering\begin{tikzpicture}[>=Stealth]
}{\end{tikzpicture}\end{beamercolorbox}}
\title{Le protocole FTP}
\subtitle{BTS CIEL 1 · Réseaux · 09 — FTP}
\author{} \date{}
\begin{document}
\begin{frame}[plain]
\vfill\etiquette{Échanger des fichiers à distance}\par\bigskip
{\Huge\bfseries Le protocole FTP}\par\bigskip
{\Large BTS CIEL 1 · Réseaux · 09 — FTP}\par\bigskip
\begin{schema}
\node[boite,font=\normalsize] (c) {Client FTP};
\node[boite,font=\normalsize,right=4cm of c] (s) {Serveur FTP};
\draw[flux] (c)--node[above]{fichier} (s);
\end{schema}\vfill
\end{frame}
\begin{frame}{Un fichier à distance}
\question{Comment récupérer un rapport sur une autre machine ?}
\begin{schema}\begin{scope}[scale=1.25,transform shape]
\node[boite,font=\normalsize] (c) {Votre ordinateur};
\node[boite,font=\normalsize,right=3cm of c] (s) {Machine distante\\rapport};
\draw[flux] (s)--node[above]{réseau} (c);
\node[below=8mm of c,align=center] {Une clé USB\\ne traverse pas le réseau.};
\end{scope}\end{schema}
\end{frame}
\begin{frame}{Un casier à code}
\question{Comment accéder au bon fichier ?}
\begin{schema}\begin{scope}[scale=1.05,transform shape]
\node[boite,font=\normalsize] (a) at (0,0) {Vous};
\node[boite,font=\normalsize] (b) at (9,0) {Casier\\avec le rapport};
\draw[flux] ([yshift=3mm]a.east)--node[above=2mm]{présenter le code} ([yshift=3mm]b.west);
\draw[flux] ([yshift=-3mm]b.west)--node[below=2mm]{prendre le rapport} ([yshift=-3mm]a.east);
\node[align=center,font=\normalsize] at (4.5,-1.8) {Code : identifiant et mot de passe\\Casier : serveur FTP};
\end{scope}\end{schema}
FTP organise l'accès et le transfert par le réseau.
\end{frame}
\begin{frame}{Exemple : récupérer un rapport}
\etiquette{Exemple fictif}\par
\begin{schema}\begin{scope}[scale=1.15,transform shape]
\node[boite,font=\normalsize,text width=3cm] at (0,0) {Client FTP\\192.168.10.20};
\node[boite,font=\normalsize,text width=3cm] at (7,0) {Serveur FTP\\192.168.10.30};
\draw[thick] (0,-0.6)--(0,-4.1); \draw[thick] (7,-0.6)--(7,-4.1);
\draw[flux] (0,-1.1)--node[above,fill=creme]{1. connexion TCP, port 21} (7,-1.1);
\draw[flux] (0,-1.8)--node[above,fill=creme]{2. utilisateur et mot de passe} (7,-1.8);
\draw[flux,dashed] (7,-2.5)--node[above,fill=creme]{accès accepté} (0,-2.5);
\draw[flux] (0,-3.2)--node[above,fill=creme]{3. demander rapport.txt} (7,-3.2);
\draw[flux] (7,-3.9)--node[above,fill=creme]{4. fichier rapport.txt} (0,-3.9);
\end{scope}\end{schema}\end{frame}
\begin{frame}{Deux connexions TCP}
\begin{schema}\begin{scope}[scale=1.05,transform shape]
\node[boite,font=\normalsize] (c) at (0,0) {Client FTP};
\node[boite,font=\normalsize] (s) at (9,0) {Serveur FTP};
\draw[flux] ([yshift=3mm]c.east)--node[above=2mm]{commandes : port serveur 21} ([yshift=3mm]s.west);
\draw[flux] ([yshift=-3mm]s.west)--node[below=2mm]{données : port serveur 20} ([yshift=-3mm]c.east);
\node[align=center,font=\normalsize] at (4.5,-1.7) {Mode actif : le serveur ouvre\\la connexion de données.};
\end{scope}\end{schema}
La connexion de commandes reste ouverte. Une liste ou un transfert ouvre une connexion de données.\\
Mode passif : le client ouvre la connexion de données vers un port indiqué par le serveur. Beaucoup de clients l'utilisent.
\end{frame}
\begin{frame}{Quel sens pour le fichier ?}
\begin{schema}\begin{scope}[scale=1.2,transform shape]
\node[boite,font=\normalsize] (c) at (0,0) {Client FTP};
\node[boite,font=\normalsize] (s) at (7,0) {Serveur FTP};
\draw[flux] ([yshift=4mm]s.west)--node[above]{télécharger} ([yshift=4mm]c.east);
\draw[flux] ([yshift=-4mm]c.east)--node[below]{téléverser} ([yshift=-4mm]s.west);
\node[font=\normalsize] at (3.5,-1.8) {Le sens se lit depuis le client.};
\end{scope}\end{schema}\end{frame}
\begin{frame}{Attention au mot de passe}
\begin{schema}\begin{scope}[scale=1.05,transform shape]
\node[boite,font=\normalsize] (c) at (0,0) {Client FTP};
\node[boite,font=\normalsize] (s) at (9,0) {Serveur FTP};
\draw[flux] (c)--node[above=2mm]{identifiant et mot de passe} (s);
\node[boite,font=\normalsize,text width=4cm] at (4.5,-1.8) {Capture : mot de passe\\soleil42 (lisible)};
\end{scope}\end{schema}
FTP transmet les identifiants en clair. Préférez SFTP (SSH : protocole d’accès à distance chiffré, port 22) ou FTPS (TLS).
\end{frame}
\begin{frame}{Méthode : réussir un transfert}
\etiquette{Invite de commandes d'un PC Packet Tracer ou de Windows}
\begin{schema}\begin{scope}[scale=1.05,transform shape]
\node[boite,text width=2.2cm] (a) at (0,0) {1. Réseau};
\node[boite,text width=2.2cm] (b) at (3,0) {2. Connexion};
\node[boite,text width=2.2cm] (c) at (6,0) {3. Liste};
\node[boite,text width=2.2cm] (d) at (9,0) {4. Transfert};
\draw[flux] (a)--(b); \draw[flux] (b)--(c); \draw[flux] (c)--(d);
\node[align=center] at (4.5,-1.1) {Vérifier \quad $\rightarrow$ \quad s'identifier \quad $\rightarrow$ \quad choisir le fichier \quad $\rightarrow$ \quad transférer};
\end{scope}\end{schema}
\begin{block}{Méthode}
\begin{enumerate}\setlength\itemsep{0pt}
\item \texttt{ping <adresse du serveur>}.
\item \texttt{ftp <adresse du serveur>}, puis identifiant et mot de passe.
\item \texttt{dir} pour lister les fichiers.
\item \texttt{get} pour télécharger, \texttt{put} pour téléverser.
\end{enumerate}
\end{block}\end{frame}
\begin{frame}{À vous : choisir le sens}
\question{Pour chaque fichier : quel sens, quel mot et quelle commande ?}
\begin{schema}\begin{scope}[scale=1.05,transform shape]
\node[boite,font=\normalsize] (c) at (0,0) {Votre PC};
\node[boite,font=\normalsize] (s) at (9,0) {Serveur\\de l'entreprise};
\draw[thick,dashed] (c)--(s);
\node[align=center,font=\normalsize] at (4.5,-1.5) {A. Mettre en ligne \texttt{accueil.html} du PC vers le serveur Web.\\[3mm] B. Récupérer sur le PC \texttt{bilan.pdf} déposé sur le serveur.};
\end{scope}\end{schema}\end{frame}
\begin{frame}{Correction : choisir le sens}
\begin{schema}\begin{scope}[scale=1.05,transform shape]
\node[boite,font=\normalsize] (c) at (0,0) {Votre PC};
\node[boite,font=\normalsize] (s) at (9,0) {Serveur\\de l'entreprise};
\draw[flux] ([yshift=3mm]c.east)--node[above=2mm]{\texttt{accueil.html} : téléverser, \texttt{put}} ([yshift=3mm]s.west);
\draw[flux] ([yshift=-3mm]s.west)--node[below=2mm]{\texttt{bilan.pdf} : télécharger, \texttt{get}} ([yshift=-3mm]c.east);
\node[align=center,font=\normalsize] at (4.5,-1.8) {A. PC vers serveur Web.\\[3mm] B. Serveur vers PC.};
\end{scope}\end{schema}\end{frame}
\begin{frame}{À vous : diagnostiquer}
\question{Le port 21 fonctionne, mais la liste échoue. Pourquoi ?}
\begin{schema}\begin{scope}[scale=1.05,transform shape]
\node[boite,font=\normalsize] (c) at (0,0) {Client FTP};
\node[boite,font=\normalsize,] (s) at (9,0) {Serveur FTP};
\draw[flux] ([yshift=3mm]c.east)--node[above=2mm]{1. commandes : OK} ([yshift=3mm]s.west);
\draw[flux,dashed] ([yshift=-4mm]s.west)--node[below=2mm]{2. liste : échec} ([yshift=-4mm]c.east);
\node[font=\normalsize] at (4.5,-1.5) {Deux connexions distinctes};
\end{scope}\end{schema}\end{frame}
\begin{frame}{Correction : diagnostiquer}
\begin{schema}\begin{scope}[scale=1.05,transform shape]
\node[boite,font=\normalsize] (c) at (0,0) {Client FTP};
\node[boite,font=\normalsize,] (s) at (9,0) {Serveur FTP};
\draw[flux] ([yshift=3mm]c.east)--node[above=2mm]{1. commandes : OK} ([yshift=3mm]s.west);
\draw[flux,dashed] ([yshift=-4mm]s.west)--node[below=2mm]{2. données : échec} ([yshift=-4mm]c.east);
\node[font=\normalsize] at (4.5,-1.5) {Deux connexions distinctes};
\end{scope}\end{schema}
La liste utilise une deuxième connexion : la connexion de données, distincte de celle du port 21. Elle échoue ici. Un pare-feu peut laisser passer le port 21 et bloquer cette connexion. Un pare-feu est un équipement ou logiciel qui autorise ou bloque des connexions selon des règles.
\end{frame}
\begin{frame}{À retenir}
\begin{schema}\begin{scope}[scale=1.05,transform shape]
\node[boite,font=\normalsize] (c) at (0,0) {Client FTP};
\node[boite,font=\normalsize,] (s) at (9,0) {Serveur FTP};
\draw[flux] ([yshift=3mm]c.east)--node[above=2mm]{commandes : TCP 21} ([yshift=3mm]s.west);
\draw[flux] ([yshift=-4mm]s.west)--node[below=2mm]{données : port 20 en mode actif} ([yshift=-4mm]c.east);
\node[font=\normalsize] at (4.5,-1.5) {Deux connexions distinctes};
\end{scope}\end{schema}
Télécharger : serveur vers client ; téléverser : client vers serveur.\\
FTP est en clair : utilisez SFTP ou FTPS.
\end{frame}
\begin{frame}{Sources et pour aller plus loin}
\begin{beamercolorbox}[sep=12pt,rounded=true]{block body}
\begin{itemize}
\item Support de l’enseignant : « TP DHCP, HTTP, DNS, FTP », p. 1--3 et 11--17.
\item RFC 959 : FTP et mode passif\\\texttt{https://www.rfc-editor.org/rfc/rfc959}
\item RFC 4217 : FTP avec TLS\\\texttt{https://www.rfc-editor.org/rfc/rfc4217}
\item RFC 4253 : SSH\\\texttt{https://www.rfc-editor.org/rfc/rfc4253}
\end{itemize}\end{beamercolorbox}
\end{frame}
\end{document}
